\section{実装概要}
\label{sec:implementation}

提案手法はUnity 6（URP）プロジェクト \texttt{HRI-Robot} に実装済みである．
\Cref{tab:impl}に論文概念と実装クラスの対応を示す．

\begin{table}[htbp]
  \centering
  \caption{論文概念と実装の対応}
  \label{tab:impl}
  \begin{tabular}{ll}
    \toprule
    論文概念 & 実装クラス \\
    \midrule
    動的基準線 & \texttt{DynamicCrossingLineTracker} \\
    危険度スコア & \texttt{VehicleRiskCalculator} \\
    色・不透明度 & \texttt{RiskToVisualMapper} \\
    経路描画 & \texttt{PathRenderer} \\
    条件切替 & \texttt{AGVPathVisualizer} \\
    実験制御 & \texttt{ExperimentManager} \\
    データ出力 & \texttt{MeasurementHub}, \texttt{ExperimentWorkbook} \\
    \bottomrule
  \end{tabular}
\end{table}

\subsection{AGVシミュレーション}

各AGVは \texttt{MovingToPickup} $\rightarrow$ \texttt{DwellAtPickup}
$\rightarrow$ \texttt{MovingToDrop} $\rightarrow$ \texttt{DwellAtDrop}
のサイクルを反復する．
経路計画はNavMesh上のA*相当処理（\texttt{AGVRoutePlanner}）で，
床面高 $y=\SI{0.15}{m}$ に固定する．
全AGVは\SI{0.1}{s}固定クロックで一括更新される（\texttt{AGVFleetOrchestrator}）．

\subsection{実験フロー}

\texttt{ExperimentStartMenu} により条件（Baseline / No-AR / Proposed）と
ケース1--10を選択する．
ケース開始時に \texttt{AGVSpawner.RestartWithSeed(seed)} を呼び，
ゴール到着時にExcelへ結果を追記する．

\subsection{再現性}

同一リポジトリ，同一シーン，同一シード列により，
他研究者が実験刺激を再現可能である．
倉庫3Dアセット（Unity Warehouse）はローカルインポートとし，
GitHubには含めない（ライセンス・容量のため）．
